\chapter{System Architecture}

\section{From Mathematics to Software}

Part IV gave this book's core cycle, observe, propose, project, commit, a complete mathematical treatment: a configuration manifold with a metric, an energy functional, a constraint manifold defined as an intersection of admissibility conditions, a sheaf-theoretic account of global consistency, conservation laws, and a way of comparing worlds to one another. None of this, on its own, runs. This Part asks what has to exist, as actual software, for the cycle Chapters 12 through 15 described to execute against a real, evolving world rather than remain a description of what a correct execution would look like. The translation from operator to component is not mechanical. Several of the boundaries this chapter draws are architectural decisions in their own right, chosen to preserve discipline this book has argued for since Part III, not merely convenient ways of organizing code.

\section{The Core Loop as a Scheduler}

Chapter 17 established that the field's components do not all update at the same rate: appearance and flow are fast, coherence and topological structure are slow, and residue inherits whatever rate the events it accumulates actually occur at. A real implementation cannot treat this as a mathematical nicety to be waved at; it has to be an actual execution policy. The component responsible for this is the scheduler, and its job is narrow but essential: at each tick of whatever clock the fast components run on, dispatch an update cycle, propose, project, commit, to the fast-timescale components at the region currently under consideration; at a coarser interval, dispatch the same cycle to the slow-timescale components. Nothing about the write cycle itself changes between these two dispatches, exactly as Chapter 17 insisted; only the rate at which the scheduler chooses to invoke it differs. A naive implementation that runs every component's update cycle on a single global tick will either waste enormous computation re-checking slow components that have not meaningfully changed, or, worse, let fast components starve while waiting for an update cycle sized for coherence and topology rather than for water moving through a fountain. The scheduler exists specifically to prevent this, and its correctness is measured against Chapter 17's timescale separation directly.

\section{The World Database}

\(M_t\) has to live somewhere queryable, and the component responsible for holding it, this chapter will call the world database, without yet specifying how it actually stores anything internally; that is Chapter 26's task, kept deliberately separate from this one. What matters here is the world database's role rather than its implementation: it is the single, authoritative location holding the field's current committed state, decomposed per Chapter 7 into Φ, v, S, R, and z at every point of whatever cover the current implementation uses, consistent with Chapter 21's patch structure. Every other component in this chapter's architecture either reads from it or, through a narrow and specific interface, writes to it, and no component is permitted to hold a private, divergent copy of any part of \(M_t\) that other components cannot see. This single-authoritative-copy discipline is what makes Chapter 15's account of commit meaningful in software terms: a commit is not complete until it has been written to this one place, and nothing counts as part of the field's history until the world database says so.

\section{The Proposal and Projection Services}

Chapters 13 and 14 kept proposal and projection as conceptually separate operators for a specific reason, argued at length in section 13.4: a proposal generator that has fully internalized every admissibility constraint loses the expressivity that makes it worth having, and a projection stage that trusts proposals unconditionally reintroduces exactly the drift Chapter 12 warned against. This chapter turns that conceptual separation into an architectural one. The proposal service, whatever neural machinery actually implements \(F_\psi\) and Chapter 13's relevance-based selection, is given read access to the world database and to whatever cue and affordance information Chapter 10's construction requires, and produces candidate \(\Delta\psi\) values. Critically, it is given no write access to the world database at all. It cannot commit anything, however confident its output. The projection service receives candidates from the proposal service, has its own read access to the world database and to whatever representation of C the implementation maintains, and computes \(\Pi_C\) against them; it, in turn, has no authority to commit either, only to produce a projected Y. Only a third component, responsible for Chapter 15's arbitration among simultaneous proposals and for the actual commit operation, has write access to the world database. This three-way separation is not merely tidy. It means the expressivity-without-authority discipline Chapter 13 argued for in the abstract is enforced by the software's own permission structure, not merely by the training or good behavior of whatever model implements the proposal service. A proposal service that hallucinates wildly is, by construction, incapable of writing that hallucination directly into the world; it can only ever offer a candidate for something else to check.

\section{The Renderer as a Client, Not a Component}

It is tempting to treat rendering, Chapter 9's concrete instance of \(O_\theta\), as another stage in the same pipeline as proposal and projection, sitting downstream of commit the way commit sits downstream of projection. This chapter argues against that framing. A renderer, in this architecture, is a client of the world database, not a stage of the core loop: it issues reads, using whatever query the observation operator requires, and produces an observation for whatever observer it serves, but it neither participates in nor is required by the write cycle itself. This distinction matters once Chapter 6's world-observer-generator separation is taken seriously as a design constraint rather than only a philosophical one. Multiple renderers, serving multiple observers, human players, other AI agents, diagnostic tools, can attach to the same world database simultaneously, each issuing its own reads, without any of them needing to coordinate with the core proposal-projection-commit loop or with one another. The world continues to evolve, through its own write cycle, whether or not any renderer happens to be attached to it at a given moment, which is precisely the property Chapter 6 insisted a genuine world needed to have.

\section{A Reference Architecture}

It is worth assembling these pieces into a single picture before the next several chapters each zoom into one part of it. At the center sits the world database, holding \(M_t\). A scheduler dispatches update cycles to it at the appropriate rate for each of the field's fast and slow components. Each cycle begins with the proposal service, reading the world database and producing candidate \(\Delta\psi\) values; continues through the projection service, reading the world database and the constraint representation to compute Y; and completes at the commit and residue ledger, which arbitrates among simultaneous proposals per Chapter 15, writes accepted results back to the world database, and updates accumulated residue. Independently of this loop, any number of renderers attach to the world database as clients, issuing reads and producing observations for whatever observers they serve, entirely decoupled from the write cycle's own rhythm. Chapter 29's multiplayer treatment will extend this picture to multiple simultaneous proposal streams arriving from different agents; Chapter 30's scaling treatment will extend it to a world database too large to hold on a single machine. Neither extension changes the shape drawn here. Both extend it.

\section{Looking Ahead}

This chapter has treated the world database as a component with a clear role and no specified internals: something that holds \(M_t\), accepts reads from proposal services and renderers, and accepts writes only from the commit stage. Chapter 26 opens that box, asking what actually holding a persistent, spatially extended, five-component semantic field requires, in terms of storage format, compression, and the practical realities of a structure too large to keep entirely in memory at once.
